\NativeSlot{45}{36}{%
% p035_s01

\ExamQuestion{4.}{5}%
判断：直角三角形和钝角三角形都只有1条高。\mbox{（\hspace{1.8em}）}\par
}
\NativeSlot{71}{68}{%
\ExamGroup{四、操作。{\normalfont\kaishu\fontsize{10.8bp}{14bp}\selectfont （共18分）}}

\par\noindent%
\ExamTile{.485}{%
% p035_s02

\ExamQuestion{1（平行四边形）}{5}%
画出下面平行四边形底边上的高。（4分）\par
\BookFigure{figures/p035_s02-01.tex}
}\hfill%
\ExamTile{.485}{%
% p035_s03

\ExamQuestion{1（三角形）}{5}%
画出下面三角形底边上的高。（4分）\par
\BookFigure{figures/p035_s03-01.tex}
}\par
}
\NativeSlot{185.0}{156}{%
% p035_s04

\ExamQuestion{2.}{5}%
在下面的点子图中画一画。（每相邻的两个点之间的距离是1 cm）\par\noindent （1）画一个底长3 cm、高2 cm的平行四边形。\par

% p035_s05

\ExamQuestion{（2）}{5}%
画一个高4 cm的等腰梯形。\par

% p035_s06

\ExamQuestion{（3）}{5}%
画一个底和高都是3 cm的钝角三角形。（9分）\par
\BookFigure{figures/p035_s04-01.tex}
}
\NativeSlot{467.5}{72}{%
% p035_s07

\ExamQuestion{3.}{5}%
已知线段\(AB=3\text{ cm}\)，请用直尺和圆规以\(AB\)为边，画一个等边三角形。（保留作图痕迹）（5分）\par
\BookFigure{figures/p035_s07-01.tex}
}
\NativeSlot{597}{115}{%
\ExamGroup{五、解决问题。{\normalfont\kaishu\fontsize{10.8bp}{14bp}\selectfont （共18分）}}

% p035_s08

\ExamQuestion{1.}{7}%
小明在打扫教室卫生时，不小心打碎了班级里的三角形玻璃教具，如图所示。你能帮他算出破碎角的度数吗？（6分）\par
\BookFigure{figures/p035_s08-01.tex}
}
